\begin{katapengantar}
    Segala puji syukur peneliti panjatkan kepada Tuhan Yesus Kristus atas kasih, berkat, dan penyertaan-Nya sehingga peneliti dapat menyelesaikan skripsi yang berjudul \textbf{"\judulindonesiatext"}. Tujuan penulisan skripsi ini adalah sebagai salah satu syarat untuk memperoleh gelar Sarjana Manajemen pada Program Studi S1 Manajemen, Fakultas Ekonomi dan Bisnis, Universitas Sumatera Utara.

    Pada kesempatan ini, peneliti mengucapkan terima kasih yang sebesar-besarnya kepada kedua \textbf{orang tua tercinta} yang telah membesarkan, mendidik, mendoakan, dan memberikan dukungan kepada peneliti. Selama proses penyusunan skripsi ini, peneliti menyadari banyak mendapat bantuan, bimbingan, dan dukungan dari berbagai pihak. Untuk itu, peneliti mengucapkan terima kasih sebesar-besarnya kepada:

    \begin{enumerate}
        \item Bapak Prof. Dr. Muryanto Amin, S.Sos., M.Si., selaku Rektor Universitas Sumatera Utara.
        \item Bapak Dr. Abdillah Arif Nasution, S.E., M.Si., Ak., C.A., QGIA, CHRS, selaku Dekan Fakultas Ekonomi dan Bisnis Universitas Sumatera Utara.
        \item Ibu Prof. Dr. Khaira Amalia Fachrudin, S.E., Ak., M.B.A., selaku Ketua Program Studi S1 Manajemen dan Bapak Muhammad Arif Lubis, S.E., M.Si., selaku Sekretaris Program Studi S1 Manajemen Fakultas Ekonomi dan Bisnis Universitas Sumatera Utara.
        \item Ibu Dr. Tetty Yuliaty, S.E., M.Si., selaku Dosen Pembimbing yang telah memberikan bimbingan, arahan, masukan, dan dukungan kepada peneliti dalam penyusunan skripsi ini.
        \item Ibu Prof. Dr. Arlina Nurbaity Lubis, S.E., M.B.A., selaku Dosen Penguji I yang telah memberikan masukan dan saran untuk penyempurnaan skripsi ini.
        \item Bapak Haryaji Catur Putera Hasman, S.E., M.Si., selaku Dosen Penguji II yang telah memberikan masukan dan saran untuk penyempurnaan skripsi ini.
        \item Seluruh dosen dan pegawai Fakultas Ekonomi dan Bisnis Universitas Sumatera Utara, khususnya Program Studi S1 Manajemen, yang telah memberikan ilmu, bantuan, dan pelayanan akademik kepada peneliti selama masa perkuliahan.
        \item Pemilik dan seluruh karyawan Rumah Makan Nasi Gerilya yang telah memberikan izin, waktu, informasi, dan bantuan selama proses penelitian.
        \item Adrian, selaku \textit{significant other}, yang selama 24/7 memberikan dukungan sepenuhnya, menemani, dan menguatkan peneliti sejak awal hingga akhir proses penyusunan skripsi ini.
        \item Sahabat peneliti, yaitu Nomi, Alta, Inggrid, Mea, Isma, dan Simon, rekan-rekan mahasiswa Program Studi S1 Manajemen, serta semua pihak yang tidak dapat disebutkan satu per satu yang telah memberikan doa, bantuan, dan dukungan kepada peneliti.
    \end{enumerate}

    Peneliti menyadari bahwa skripsi ini masih jauh dari sempurna. Oleh karena itu, peneliti mengharapkan kritik dan saran yang membangun demi penyempurnaan skripsi ini. Akhir kata, peneliti berharap skripsi ini dapat memberikan manfaat bagi pembaca dan pihak-pihak yang membutuhkan.
\end{katapengantar}
